% Archivo autónomo en UTF-8. No necesita imágenes ni archivos externos.
% Formato APA 7 para trabajo estudiantil. Completar los datos de portada.
% Compilar dos veces con pdfLaTeX para actualizar referencias.
% También admite LuaLaTeX o XeLaTeX. No requiere shell-escape.
\documentclass[12pt,letterpaper]{article}
\usepackage{iftex}
\ifPDFTeX
  \usepackage[utf8]{inputenc}
  \usepackage[T1]{fontenc}
  \usepackage{newtxtext}
\else
  \usepackage{fontspec}
  \setmainfont{TeX Gyre Termes}
  \setsansfont{Latin Modern Sans}
  \setmonofont{Latin Modern Mono}
\fi
\usepackage[spanish,es-tabla,es-nodecimaldot]{babel}
\usepackage[margin=1in]{geometry}
\usepackage{microtype}
\usepackage{amsmath}
\usepackage{array,booktabs,longtable,tabularx}
\usepackage{enumitem}
\usepackage{xcolor}
\usepackage{xurl}
\usepackage{tikz}
\usetikzlibrary{arrows.meta,positioning,babel}
\usepackage{hyperref}
\definecolor{quickly}{HTML}{BE185D}
\definecolor{ink}{HTML}{172033}
\definecolor{soft}{HTML}{FDF2F8}
\hypersetup{hidelinks,
  pdftitle={Quickly Delivery: manual de funcionamiento},
  pdfauthor={Proyecto Frontend-Quickly},unicode=true}
\setlength{\parindent}{0.5in}
\setlength{\parskip}{0pt}
\setlength{\emergencystretch}{2em}
\setlist{itemsep=0.25em,topsep=0.3em}
\renewcommand{\arraystretch}{1.0}
\newcolumntype{L}[1]{>{\raggedright\arraybackslash}p{#1}}
\newcommand{\codigo}[1]{\texttt{\detokenize{#1}}}
\newcommand{\soles}[1]{S/\,#1}

\usepackage{setspace}
\usepackage{fancyhdr}
\usepackage{titlesec}
\usepackage{caption}
\usepackage{ragged2e}
\doublespacing
\setcounter{secnumdepth}{0}
\pagestyle{fancy}
\fancyhf{}
\fancyhead[R]{\thepage}
\renewcommand{\headrulewidth}{0pt}
\setlength{\headheight}{15pt}
\setlength{\LTcapwidth}{\textwidth}
\titleformat{\section}{\normalfont\bfseries\centering}{}{0pt}{}
\titleformat{\subsection}{\normalfont\bfseries}{}{0pt}{}
\titleformat{\subsubsection}{\normalfont\bfseries\itshape}{}{0pt}{}
\captionsetup{position=top,labelfont=bf,textfont=it,labelsep=newline,
  justification=raggedright,singlelinecheck=false,font=normalsize}
% Completar los datos reales antes de entregar el trabajo.
\newcommand{\nombreAutor}{[Nombre y apellido del autor]}
\newcommand{\institucion}{[Departamento e institución]}
\newcommand{\asignatura}{[Código y nombre de la asignatura]}
\newcommand{\docente}{[Nombre del docente]}
\newcommand{\fechaEntrega}{[Fecha de entrega]}
\newcommand{\apaportada}{%
  \thispagestyle{fancy}
  \vspace*{1in}
  \begin{center}
    \textbf{Quickly Delivery: manual de funcionamiento del sistema}\par
    \vspace{\baselineskip}
    \nombreAutor\par
    \institucion\par
    \asignatura\par
    \docente\par
    \fechaEntrega\par
  \end{center}
}

\begin{document}
\apaportada
\clearpage
\renewcommand{\contentsname}{Índice}
\setcounter{tocdepth}{2}
\tableofcontents
\clearpage
\RaggedRight
\setlength{\parindent}{0.5in}
\begin{center}\textbf{Quickly Delivery: manual de funcionamiento del sistema}\end{center}


\section{Propósito y alcance}

Quickly Delivery permite explorar comercios, seleccionar productos, registrar una compra y coordinar su preparación y entrega. Sus interfaces representan cuatro participantes: cliente, comercio, repartidor y administrador.

La presentación adopta el formato de trabajo estudiantil de APA 7 (American Psychological Association, s. f.).

Este manual describe el funcionamiento implementado en el repositorio, con especial atención al flujo de Carrito y Checkout y a los paneles de administración y reparto. Sirve como guía de uso y como referencia para mantenimiento. Las rutas y archivos corresponden al commit \codigo{d0a8c92} del repositorio (xmenil, 2026); cambios posteriores pueden modificar este comportamiento.

\subsection{Estado actual del proyecto}

El sistema es un \textbf{prototipo frontend interactivo}. Las acciones cambian datos compartidos dentro de la aplicación y los guardan en el navegador. No existe en esta versión un servidor central de pedidos, una pasarela de pagos conectada ni seguimiento GPS real.

Los pedidos y direcciones precargados permiten observar pantallas con contenido y probar operaciones. Su presencia no demuestra compras, cobros o entregas reales. El registro local de un estado de pago tampoco constituye una confirmación bancaria.

\subsection{Conceptos principales}

\begin{description}
  \item[Compra.] Registro general que reúne cliente, dirección, método de pago, importes y pedidos por comercio. En el código se representa mediante \codigo{Purchase}.
  \item[Pedido por comercio.] Parte de una compra asociada a un local, con productos, estado, repartidor y eventos propios. Se representa mediante \codigo{MerchantOrder}.
  \item[Carrito.] Productos elegidos por el cliente, incluidos los que conserva para después. Solamente los seleccionados se compran.
  \item[Zona de cobertura.] Sector con tarifa base, tiempo estimado y disponibilidad de reparto.
  \item[Incidencia.] Caso registrado para revisión de soporte o administración.
\end{description}

La interfaz actual permite comprar de \textbf{un comercio a la vez}. El modelo conserva un arreglo de pedidos por comercio para representar registros anteriores y estructuras multicomercio. Esta capacidad del modelo no implica que el carrito actual permita pagar varios locales simultáneamente.

\section{Participantes y módulos}

\begin{table}[htbp]
\caption{Participantes y responsabilidades}
\begin{singlespace}
\begin{tabularx}{\textwidth}{L{2.8cm}X}
\toprule
Participante & Responsabilidad principal \\
\midrule
Cliente & Seleccionar productos, indicar dirección y pago, registrar el pedido y consultar su estado. \\
Comercio & Administrar su catálogo y disponibilidad, revisar pedidos y avanzar su preparación. \\
Repartidor & Aceptar entregas disponibles, confirmar recojo y entrega, y reportar inconvenientes. \\
Administrador & Revisar solicitudes, supervisar cuentas, comercios, repartidores, pedidos e incidencias. \\
\bottomrule
\end{tabularx}
\end{singlespace}
\noindent\textit{Nota.} Elaboración propia a partir del código del proyecto (xmenil, 2026).
\end{table}

\begin{figure}[htbp]
\caption{Relación operativa entre los paneles. En el prototipo utilizan datos locales compartidos.}
\centering
\begin{tikzpicture}[node distance=0.45cm,
  panel/.style={draw=quickly,fill=soft,rounded corners,align=center,
    text width=2.7cm,minimum height=1.1cm,font=\small\sffamily},
  arrow/.style={-{Latex},thick,draw=ink}]
\node[panel] (client) {Cliente\\Registra la compra};
\node[panel,right=of client] (merchant) {Comercio\\Prepara el pedido};
\node[panel,right=of merchant] (courier) {Repartidor\\Recoge y entrega};
\draw[arrow] (client) -- (merchant);
\draw[arrow] (merchant) -- (courier);
\node[panel,below=0.8cm of merchant,text width=5cm] (admin) {Administrador\\Supervisa pedidos y atiende incidencias};
\draw[arrow,dashed] (merchant) -- (admin);
\draw[arrow,dashed] (courier.south) |- (admin.east);
\end{tikzpicture}
\par\medskip
\noindent\textit{Nota.} Elaboración propia basada en el flujo implementado (xmenil, 2026).
\end{figure}

\subsection{Acceso y cuentas de ejemplo}

El acceso se ofrece en \path{/login}. El prototipo contiene estas cuentas de referencia:

\begin{table}[htbp]
\caption{Cuentas de ejemplo}
\begin{singlespace}
\begin{tabularx}{\textwidth}{L{2.7cm}X}
\toprule
Rol & Correo de ejemplo \\
\midrule
Cliente & \path{cliente@quickly.pe} \\
Comercio & \path{comercio@quickly.pe} \\
Repartidor & \path{repartidor@quickly.pe} \\
Administrador & \path{admin@quickly.pe} \\
\bottomrule
\end{tabularx}
\end{singlespace}
\noindent\textit{Nota.} Elaboración propia a partir del código del proyecto (xmenil, 2026).
\end{table}

La sesión del prototipo se resuelve con los datos de ejemplo y se almacena localmente. El botón de Google crea una sesión de ejemplo; no ejecuta una integración OAuth real. Estas cuentas no son credenciales de un servicio de producción.

El selector de paneles permite cambiar de rol y observar el mismo conjunto de datos. Se oculta en Carrito y Checkout para no interferir con la compra. El panel de comercio también permite seleccionar el comercio activo según las funciones de la sesión local.

\subsection{Registros y aprobaciones}

El administrador puede aprobar, rechazar, suspender y reactivar usuarios existentes en el almacén de datos. Las solicitudes precargadas permiten probar ese recorrido.

Los formularios públicos de afiliación de comercio y repartidor muestran una confirmación visual, pero sus manejadores actuales no crean una nueva solicitud persistente en el almacén compartido. El registro de cliente crea una sesión local; su incorporación completa al listado administrativo también requiere integración. Por ello, no debe asumirse que todos los registros públicos aparecen automáticamente en Administración.

\section{Flujo del cliente}

\subsection{Exploración y selección de productos}

El cliente puede consultar el inicio, el catálogo de comercios, el detalle de un comercio y el detalle de un producto. Allí observa la información disponible del catálogo, precios, opciones o variantes y disponibilidad.

Agregar un producto actualiza el carrito. Si el producto corresponde a otro comercio, las reglas de selección mantienen activa una sola tienda para la compra y conservan los demás productos sin seleccionar.

\subsection{Carrito}

En \path{/carrito} los productos se agrupan por comercio. Cada grupo presenta el nombre del local, su disponibilidad, tiempo de preparación publicado y pedido mínimo. El cliente puede:

\begin{itemize}
  \item Elegir ``Comprar de este comercio'' para seleccionar sus productos.
  \item Marcar o desmarcar productos mediante casillas explícitas.
  \item Aumentar o disminuir cantidades dentro de la disponibilidad permitida.
  \item Eliminar un producto con su acción específica.
  \item Vaciar todo el carrito después de confirmar la eliminación, incluidos los productos guardados para después.
\end{itemize}

Disminuir una cantidad de una unidad está bloqueado en la interfaz: para retirarla se utiliza la acción de eliminar. Modificar cantidades no cambia por sí mismo la selección de comercio.

El resumen calcula únicamente lo seleccionado. Antes de elegir dirección, el envío y el total se presentan como \textbf{estimados}. Si no existe una zona elegida, la función de cálculo utiliza la tarifa del comercio o su valor de respaldo. En Checkout se utiliza la tarifa de la zona de la dirección.

La continuación se bloquea si faltan productos seleccionados, el comercio no está activo o abierto, un producto u opción no está disponible, el stock no alcanza, el precio está desactualizado o no se llega al pedido mínimo. Si un precio cambió, la interfaz indica retirar el producto y agregarlo nuevamente.

Si el cliente no ha iniciado sesión, continuar conduce al acceso y conserva la intención de regresar a Checkout.

\subsection{Checkout y protección de pasos}

El recorrido de \path{/checkout} es:

\begin{center}
\textbf{Dirección} $\longrightarrow$ \textbf{Pago} $\longrightarrow$ \textbf{Revisión} $\longrightarrow$ \textbf{Pedido registrado}
\end{center}

\begin{enumerate}
  \item \textbf{Dirección:} seleccionar una dirección de la cuenta, con cobertura disponible, y continuar al pago.
  \item \textbf{Pago:} elegir expresamente un método y completar sus condiciones. No existe un método seleccionado por defecto.
  \item \textbf{Revisión:} comprobar productos, dirección, tarifa y pago; corregir los datos mediante sus acciones y confirmar el importe final.
\end{enumerate}

El indicador permite volver a pasos anteriores o a pasos ya completados. No habilita la revisión con solo haber seleccionado una dirección. Cambiar los productos, la dirección, la tarifa o datos relevantes invalida la revisión anterior. Cambiar el método de pago obliga a revisar nuevamente ese paso.

La confirmación repite las comprobaciones de sesión, selección, cobertura, comercio, precios, cantidades y stock. También comprueba si los datos cambiaron durante la espera de registro. Un bloqueo de envío evita que pulsaciones repetidas generen pedidos duplicados durante esa operación.

\subsection{Métodos de pago}

\subsubsection{Efectivo contra entrega}

El cliente indica el monto disponible. Puede utilizar coma o punto decimal, elegir el monto exacto o un atajo suficiente para cubrir el total. El sistema rechaza importes incompletos, negativos, con precisión inválida o menores que el pedido.

El vuelto se calcula como:
\[
  V = E - T, \qquad E \geq T,
\]
donde $E$ es el efectivo indicado y $T$ el total. El importe representa con cuánto \emph{pagará} el cliente, no un cobro ya recibido.

\subsubsection{Yape/Plin y tarjeta}

Elegir un método digital no genera un cargo. La interfaz informa que el pago quedará pendiente de verificación y exige reconocer esa condición antes de pasar a revisión. No ofrece un QR cobrable, un número de cobro oficial ni un formulario para recopilar datos de tarjeta.

Todas las compras nuevas se registran con pago \codigo{pendiente}. La confirmación de un cobro digital requiere una pasarela o un procedimiento administrativo que aún no está integrado. En efectivo, confirmar la entrega actualiza el estado de pago del pedido por comercio; ese cambio no garantiza que el estado global de la compra se actualice del mismo modo.

\subsection{Registro y recuperación ante errores}

Al registrar la compra, se construyen sus pedidos por comercio, se conserva una copia de productos y dirección, se reserva stock y se crean eventos y notificaciones locales. Se retiran del carrito solamente los productos adquiridos. Los demás permanecen disponibles para otra compra.

La aplicación abre el detalle del pedido y muestra ``Pedido registrado''. Si el registro falla, conserva el carrito y comunica el motivo. Si el pedido ya se guardó pero falla la limpieza del carrito, bloquea una nueva confirmación y ofrece abrir el registro existente.

\section{Direcciones, cobertura y tarifas}

\subsection{Direcciones guardadas}

En \path{/cliente/direcciones} se pueden consultar, agregar, editar, eliminar y elegir una dirección principal. Cada dirección pertenece a una cuenta y contiene destinatario, celular, calle, número, referencia, zona y nombre identificador.

La precarga temporal proporciona dos ubicaciones:

\begin{table}[htbp]
\caption{Direcciones precargadas}
\begin{singlespace}
\begin{tabularx}{\textwidth}{L{3cm}X}
\toprule
Nombre & Ubicación y referencia \\
\midrule
Casa & Av. Alameda Perú 315. Frente al Parque Ramón Castilla, puerta verde. Zona Centro. \\
Trabajo / UNAS & Av. Universitaria s/n, Pabellón Central. Facultad de Agronomía, garita 2. Zona Rupa Rupa Norte / UNAS. \\
\bottomrule
\end{tabularx}
\end{singlespace}
\noindent\textit{Nota.} Elaboración propia a partir del código del proyecto (xmenil, 2026).
\end{table}

Los datos de contacto se adaptan a la cuenta que abre las direcciones. La carga se realiza una sola vez por cuenta, preserva registros editados y no vuelve a insertar una dirección eliminada después de la precarga. Si ya existe una dirección principal, no la sustituye.

\subsection{Formulario de nueva dirección}

El modal organiza el formulario en contacto, ubicación y opciones de guardado. El celular comienza vacío, acepta separadores habituales y se guarda como nueve dígitos. Los campos de texto se limpian de espacios exteriores. El número admite referencias como manzana, lote o s/n.

La zona debe elegirse expresamente y tener cobertura al guardar. Las validaciones aparecen junto a los campos y se enfoca el primer error. Si falla el guardado, se conserva lo escrito. En Checkout, la dirección recién guardada queda seleccionada y actualiza el envío y el total, sin saltar el paso de Pago.

\subsection{Zonas de referencia}

\begin{table}[htbp]
\caption{Zonas de cobertura}
\begin{singlespace}
\begin{tabularx}{\textwidth}{XrrL{2.3cm}}
\toprule
Zona & Tarifa base & Tiempo & Disponibilidad inicial \\
\midrule
Centro de Tingo María & \soles{4.00} & 20 min & Habilitada \\
Rupa Rupa Norte / UNAS & \soles{5.00} & 25 min & Habilitada \\
Rupa Rupa Sur / Afilador & \soles{5.50} & 30 min & Habilitada \\
Castillo Grande & \soles{6.50} & 35 min & Habilitada \\
Naranjillo Periférica & \soles{8.00} & 45 min & Sin cobertura \\
\bottomrule
\end{tabularx}
\end{singlespace}
\noindent\textit{Nota.} Elaboración propia a partir del código del proyecto (xmenil, 2026).
\end{table}

Son valores iniciales del prototipo. Los tiempos son estimaciones registradas, no predicciones GPS ni compromisos de llegada. El panel administrativo de zonas permite consultarlos; no contiene actualmente un editor de tarifas.

\section{Panel del comercio}

\subsection{Organización}

\begingroup\singlespacing
\begin{longtable}{L{5.2cm}L{10.1cm}}
\caption{Rutas del comercio}\\
\toprule
Ruta & Uso principal \\
\midrule
\endhead
\path{/comercio/tienda} & Resumen operativo y disponibilidad del comercio. \\
\path{/comercio/perfil-tienda} & Datos y presentación del local. \\
\path{/comercio/productos} & Creación, edición y disponibilidad del catálogo. \\
\path{/comercio/pedidos} & Revisión de pedidos, preparación y acciones operativas. \\
\path{/comercio/reportes} & Consulta de información comercial y reportes. \\
\path{/comercio/configuracion} & Opciones de configuración del panel. \\
\bottomrule
\end{longtable}
\endgroup
\noindent\textit{Nota.} Elaboración propia a partir del código del proyecto (xmenil, 2026).\par\bigskip

\subsection{Preparación de un pedido}

El comercio consulta los pedidos asociados a su identificador. Un pedido nuevo queda pendiente hasta que se confirma o rechaza. La preparación utiliza los estados permitidos: confirmado, en preparación y listo para recojo.

Marcar un pedido listo para recojo lo incorpora a las solicitudes de reparto si todavía no tiene repartidor. Cada cambio conserva un evento en la línea de tiempo y puede generar registros de auditoría y notificaciones locales.

Pausar el comercio cambia su recepción de pedidos; no elimina el catálogo ni los pedidos existentes. El carrito y Checkout revisan su disponibilidad antes de continuar y registrar una nueva compra.

\section{Panel del repartidor}

\begingroup\singlespacing
\begin{longtable}{L{5.8cm}L{9.5cm}}
\caption{Rutas del repartidor}\\
\toprule
Ruta & Uso principal \\
\midrule
\endhead
\path{/repartidor/solicitudes} & Pedidos listos para recojo sin repartidor asignado. \\
\path{/repartidor/entrega-activa} & Dirección, contacto, cobro, recojo y entrega en curso. \\
\path{/repartidor/historial} & Entregas completadas y detalles desplegables. \\
\path{/repartidor/ingresos} & Tarifas registradas, número de entregas y promedio. \\
\path{/repartidor/perfil} & Identidad, vehículo, zona, estado de cuenta y disponibilidad. \\
\bottomrule
\end{longtable}
\endgroup
\noindent\textit{Nota.} Elaboración propia a partir del código del proyecto (xmenil, 2026).\par\bigskip

\subsection{Disponibilidad y solicitudes}

El control de cabecera permite activar o pausar la disponibilidad. Una cuenta activa y la disponibilidad son conceptos distintos: el primero corresponde a la gestión administrativa; el segundo a la disposición para recibir trabajo.

Una solicitud muestra código, comercio, tarifa, dirección de recojo, destino, referencia y método de pago. La interfaz bloquea aceptar mientras el repartidor esté en pausa o tenga una entrega en curso. Si otro repartidor tomó el pedido, la acción comunica el error.

Estos controles pertenecen a la interfaz y no sustituyen validaciones de un servidor. Algunas operaciones internas del almacén y la asignación administrativa todavía no aplican todas las restricciones de capacidad o identidad.

\subsection{Entrega activa}

La dirección y el siguiente paso aparecen primero. Durante el recojo se muestran dirección y contacto del comercio; durante la entrega, dirección, referencia y contacto del cliente. La llamada utiliza un enlace telefónico.

\begin{enumerate}
  \item Aceptar la solicitud asigna el pedido al repartidor.
  \item ``Confirmar recojo'' cambia el pedido de listo para recojo a en camino.
  \item ``Confirmar entrega'' abre una confirmación final.
  \item ``Finalizar entrega'' registra el estado entregado y abre el historial cuando la operación tiene éxito.
\end{enumerate}

Los errores conservan el contexto. El recorrido de referencia es un esquema de origen y destino, sin coordenadas ni distancia calculada. Para compras antiguas multicomercio, el vuelto global requiere cuidado: el dato de vuelto puede pertenecer a la compra completa y no al pedido individual.

\subsection{Historial, ganancias y datos precargados}

El historial muestra pedidos asignados al repartidor que estén entregados. Ganancias suma sus tarifas de entrega y calcula el promedio. Sin registros, muestra cero. La tarifa registrada no confirma una liquidación ni un pago al repartidor.

Los datos iniciales adicionales incluyen tres solicitudes listas, una entrega en camino y cuatro entregas completadas. La carga evita duplicar sus identificadores y preserva una entrega activa existente. Se integran en los datos compartidos, por lo que también pueden verse desde cliente, comercio y administración según la cuenta y el filtro utilizados.

\section{Administración y soporte}

\begingroup\singlespacing
\begin{longtable}{L{5.2cm}L{10.1cm}}
\caption{Rutas de administración}\\
\toprule
Ruta & Función \\
\midrule
\endhead
\path{/admin/resumen} & Pendientes, disponibilidad operativa y cifras generales. \\
\path{/admin/usuarios} & Solicitudes, aprobación, rechazo, suspensión y reactivación. \\
\path{/admin/comercios} & Datos de comercios y apertura o pausa. \\
\path{/admin/repartidores} & Flota, vehículo, zona y disponibilidad. \\
\path{/admin/pedidos} & Monitor de pedidos y asignación manual de repartidor. \\
\path{/admin/zonas} & Consulta de cobertura, tarifas y tiempos. \\
\path{/admin/incidencias} & Casos abiertos, en revisión y resueltos. \\
\path{/admin/reportes} & Totales registrados y exportación CSV. \\
\path{/admin/auditoria} & Historial de actividad, búsqueda y filtro por responsable. \\
\bottomrule
\end{longtable}
\endgroup
\noindent\textit{Nota.} Elaboración propia a partir del código del proyecto (xmenil, 2026).\par\bigskip

\subsection{Gestión de cuentas}

La aprobación activa al usuario y, según el caso, el perfil asociado de comercio o repartidor. El rechazo de una solicitud deja al usuario suspendido y el perfil asociado rechazado. Suspender o reactivar una cuenta utiliza acciones diferentes de pausar la disponibilidad operativa.

El rechazo y la suspensión requieren un motivo escrito en la interfaz. El sistema registra la decisión en el historial local. Estos registros son editables desde el entorno del navegador y no constituyen una auditoría inmutable de producción.

\subsection{Incidencias}

El repartidor puede informar lluvia, tráfico, falta de respuesta del cliente, dificultad para ubicar la dirección o avería. El envío crea un caso de prioridad alta, vinculado al código de compra y en revisión.

Administración consulta el caso y escribe una respuesta obligatoria para resolverlo. La interfaz conserva el texto si falla la operación y permite leer respuestas anteriores. Resolver una incidencia no modifica automáticamente el tiempo estimado, la posición del repartidor ni el estado de entrega.

\subsection{Reportes}

Los totales administrativos incluyen los registros disponibles, sin filtro de fecha y con todos sus estados. El monto total registrado incluye compras canceladas; no equivale a ingresos cobrados. Las tarifas de envío registradas tampoco equivalen a pagos liquidados a repartidores.

La exportación CSV incluye identificador, código, cliente, fecha, método de pago, total en soles y estado. El archivo procede de los datos locales del navegador.

\section{Estados y cálculos}

\subsection{Estados del pedido por comercio}

\begin{figure}[htbp]
\caption{Transiciones permitidas por las reglas del prototipo.}
\centering
\begin{tikzpicture}[node distance=0.38cm,
  state/.style={draw=quickly,fill=soft,rounded corners,text width=5cm,align=center,minimum height=0.75cm,font=\small\sffamily},
  branch/.style={draw=ink,rounded corners,text width=2.4cm,align=center,minimum height=0.75cm,font=\small\sffamily},
  arrow/.style={-{Latex},draw=ink}]
\node[state] (pending) {Pendiente};
\node[state,below=of pending] (confirmed) {Confirmado};
\node[state,below=of confirmed] (preparing) {En preparación};
\node[state,below=of preparing] (ready) {Listo para recojo};
\node[state,below=of ready] (transit) {En camino};
\node[state,below=of transit] (delivered) {Entregado};
\draw[arrow] (pending) -- (confirmed);
\draw[arrow] (confirmed) -- (preparing);
\draw[arrow] (preparing) -- (ready);
\draw[arrow] (ready) -- (transit);
\draw[arrow] (transit) -- (delivered);
\node[branch,right=1cm of pending] (rejected) {Rechazado};
\node[branch,left=1cm of confirmed] (cancelled) {Cancelado};
\draw[arrow] (pending.east) -- (rejected.west);
\draw[arrow] (pending.west) -| (cancelled.north);
\draw[arrow] (confirmed.west) -- (cancelled.east);
\end{tikzpicture}
\par\medskip
\noindent\textit{Nota.} Elaboración propia basada en el flujo implementado (xmenil, 2026).
\end{figure}

\begin{table}[htbp]
\caption{Estados del pedido}
\begin{singlespace}
\begin{tabularx}{\textwidth}{L{4.2cm}X}
\toprule
Estado interno & Significado \\
\midrule
\codigo{pendiente} & Espera aceptación del comercio. \\
\codigo{confirmado} & El comercio aceptó el pedido. \\
\codigo{en_preparacion} & El comercio prepara los productos. \\
\codigo{listo_recoger} & Puede recogerse y pasar al reparto. \\
\codigo{en_camino} & El repartidor confirmó el recojo. \\
\codigo{entregado} & Se confirmó la entrega. \\
\codigo{cancelado} & El pedido fue cancelado. \\
\codigo{rechazado} & El comercio rechazó el pedido. \\
\bottomrule
\end{tabularx}
\end{singlespace}
\noindent\textit{Nota.} Elaboración propia a partir del código del proyecto (xmenil, 2026).
\end{table}

La compra general deriva su estado de los pedidos que contiene: pendiente, en proceso, entrega parcial, completada, cancelada o finalizada con cancelaciones. Los estados globales y los estados por comercio no son intercambiables.

\subsection{Cancelaciones y reintegros}

El cliente puede cancelar pedidos pendientes o confirmados, antes de la preparación. Las transiciones terminales no permiten continuar al siguiente paso de reparto. Una cancelación o rechazo restaura stock según la implementación actual y libera la asignación del repartidor cuando corresponde.

La regla de reintegro de métodos digitales conserva una suposición anterior de prepago y puede marcar el registro como reembolsado aunque el Checkout nuevo deje el pago pendiente. \textbf{No se realiza una transferencia bancaria.} Esta regla debe reconciliarse con la verificación de pagos antes de utilizar el sistema para cobros reales.

\subsection{Importes en centavos}

El código guarda dinero como enteros en centavos y lo presenta mediante \codigo{formatCents}. Para una línea de producto:
\[
  L_i = q_i\,(p_i + v_i), \qquad
  P = \sum_i L_i, \qquad T = P + F,
\]
donde $q_i$ es la cantidad, $p_i$ el precio base, $v_i$ la diferencia de variante, $P$ el subtotal y $F$ el envío. La tarifa de la zona se aplica cuando hay dirección seleccionada. La comprobación de stock considera la cantidad conjunta del mismo producto en las opciones seleccionadas.

\subsection{Ejemplo de uso completo}

Para ilustrar el recorrido, supóngase un producto de \soles{28.00}, entrega al Centro de \soles{4.00} y pago en efectivo con \soles{50.00}. El total es \soles{32.00} y el vuelto \soles{18.00}. Son valores ilustrativos, no un cobro real.

\begin{enumerate}
  \item El cliente selecciona el comercio y el producto en Carrito.
  \item En Checkout elige Casa, comprueba el envío y continúa.
  \item Elige efectivo, indica el monto y revisa el vuelto.
  \item Revisa el pedido y confirma. Se reserva stock y se genera un código.
  \item El comercio confirma, prepara y marca listo para recojo.
  \item Un repartidor disponible, sin otra entrega activa en su interfaz, acepta la solicitud.
  \item Confirma el recojo, consulta dirección y referencia y realiza la entrega.
  \item Confirma la entrega y consulta su historial y tarifa registrada.
  \item Administración puede observar los mismos registros y atender una incidencia vinculada.
\end{enumerate}

La coincidencia entre paneles depende del mismo origen web y almacenamiento local, de la cuenta o comercio seleccionados y de los filtros. Abrir otra computadora no comparte automáticamente estos datos.

\section{Arquitectura y mantenimiento}

\subsection{Tecnologías}

La aplicación utiliza React 18 y TypeScript, Vite para desarrollo y compilación, React Router 6 para rutas, Zustand para estado, Tailwind CSS 3 para estilos, Lucide para iconos y React Hook Form con Zod para el formulario de direcciones. Las pruebas utilizan Vitest y Testing Library.

\subsection{Almacenes y persistencia}

\begin{table}[htbp]
\caption{Almacenes de datos}
\begin{singlespace}
\begin{tabularx}{\textwidth}{L{3cm}X}
\toprule
Archivo & Responsabilidad \\
\midrule
\path{src/store/authStore.ts} & Usuario, sesión, rol y comercio activo de ejemplo. \\
\path{src/store/cartStore.ts} & Productos, cantidades y selección de un comercio. \\
\path{src/store/dataStore.ts} & Compras, catálogo, direcciones, flota, incidencias, eventos y acciones compartidas. \\
\bottomrule
\end{tabularx}
\end{singlespace}
\noindent\textit{Nota.} Elaboración propia a partir del código del proyecto (xmenil, 2026).
\end{table}

Las claves principales de \codigo{localStorage} son \codigo{quickly_auth_v1}, \codigo{quickly_cart_v1} y \codigo{quickly_datastore_v4}. El almacén de datos puede recuperar versiones anteriores y agregar semillas faltantes sin reemplazar todos los registros.

Los cambios de sesión, carrito y datos cuentan con manejadores de eventos de almacenamiento para ciertas actualizaciones entre pestañas del mismo origen. Esto no equivale a sincronización con un servidor. Borrar el almacenamiento puede eliminar la sesión y las modificaciones locales.

La acción ``Restaurar datos'' repone los datos iniciales del almacén compartido. Puede borrar cambios locales de pedidos, productos e incidencias; no debe interpretarse como una copia de seguridad. Los ejemplos de reparto y direcciones cuentan con controles para evitar su inserción repetida.

\subsection{Datos y reglas}

\begin{itemize}
  \item \path{src/domain/types.ts}: entidades y tipos.
  \item \path{src/domain/pricingRules.ts}: agrupación de productos, importes y efectivo.
  \item \path{src/domain/orderRules.ts}: transiciones, estado global y cancelaciones.
  \item \path{src/features/checkout/checkoutFlow.ts}: validaciones de selección, cobertura, precios, stock y clave de revisión.
  \item \path{src/mocks/initialData.ts}: catálogo, cuentas y datos iniciales.
  \item \path{src/mocks/courierActivityData.ts}: actividad adicional de reparto.
  \item \path{src/mocks/exampleAddresses.ts}: precarga de direcciones por cuenta.
  \item \path{src/lib/currency.ts}: conversión y presentación monetaria.
  \item \path{src/lib/date.ts}: presentación de fechas en la zona horaria de Lima.
\end{itemize}

\subsection{Diseño y componentes}

Las reglas visuales se documentan en las skills locales \codigo{quickly-ui} y \codigo{ux-principles-reviewer}. Los tokens se definen en \path{src/styles/index.css} y se exponen en \path{tailwind.config.js}. Los componentes compartidos se encuentran en \path{src/components/ui} y \path{src/components/shared}.

El diseño prioriza contraste, etiquetas visibles, importes tabulares, acciones táctiles de al menos 44 píxeles, mensajes de error y conservación de datos. El diálogo compartido gestiona foco, navegación de teclado y cierre con Escape. La verificación visual en teléfonos sigue siendo necesaria; las pruebas automatizadas no comprueban por sí solas todos los tamaños o desbordamientos.

\subsection{Ejecución y comprobaciones}

Desde la raíz del proyecto:
\begin{verbatim}
npm install
npm run dev
npm test
npm run build
npm run preview
\end{verbatim}

\codigo{dev} inicia el servidor de desarrollo; \codigo{test} ejecuta las pruebas; \codigo{build} comprueba TypeScript y genera los recursos de producción en \codigo{dist}; \codigo{preview} permite consultar esa compilación localmente. \codigo{npm install} requiere Node.js y npm disponibles y acceso a las dependencias.

En la verificación previa del rediseño se aprobaron 71 pruebas y la compilación. Esa cifra describe la revisión realizada, no sustituye ejecutar las comprobaciones después de nuevos cambios. La configuración de Vercel incluye una reescritura a \codigo{index.html} para las rutas de la aplicación.

\section{Alcances pendientes y comprobación de uso}

\subsection{Integraciones pendientes}

Para evolucionar el prototipo a una operación con usuarios reales se necesita:
\begin{itemize}
  \item Servidor y base de datos central para cuentas, catálogo, compras e incidencias.
  \item Autenticación real y autorización por rol y propiedad del registro.
  \item Validaciones de stock y precios en servidor, con control de concurrencia e idempotencia.
  \item Pasarela o proceso verificable de cobro, conciliación y devolución.
  \item Registro persistente de afiliaciones y sincronización de nuevas cuentas.
  \item Reglas de asignación que impidan entregas simultáneas no autorizadas y validen quién puede avanzar un pedido.
  \item Geolocalización y cálculo de rutas si se desea seguimiento real.
  \item Notificaciones externas y una auditoría protegida frente a modificación.
\end{itemize}

Actualmente, los layouts de comercio, repartidor y administrador no deben considerarse barreras de autorización de producción. Algunos seleccionan el primer perfil disponible como respaldo. Los controles de Checkout y de selección de reparto mejoran el flujo normal de la interfaz, pero no impiden manipular el estado desde el navegador.

\subsection{Recorrido recomendado para probar el sistema}

\begin{enumerate}
  \item Iniciar sesión con el cliente de ejemplo y agregar un producto de un comercio abierto.
  \item Confirmar que cambiar de comercio conserva los otros productos sin seleccionar.
  \item Verificar que no se puede abrir Revisión antes de Dirección y Pago.
  \item Agregar una dirección y comprobar su selección y el cambio de tarifa.
  \item Probar efectivo insuficiente, monto exacto y vuelto.
  \item Registrar un pago digital y comprobar que permanece pendiente.
  \item Confirmar una compra y comprobar que solo desaparece lo adquirido del carrito.
  \item Cambiar al comercio correspondiente y avanzar los estados permitidos.
  \item Cambiar al repartidor, completar su entrega activa si la tiene y aceptar un pedido listo.
  \item Reportar una incidencia y resolverla desde Administración.
  \item Confirmar entrega y consultar Historial, Ganancias y Auditoría.
\end{enumerate}

También deben probarse los estados vacíos, zonas sin cobertura, cambios de precio o stock durante la compra, cancelaciones permitidas y fallos de almacenamiento. La revisión visual debe incluir escritorio y teléfono, navegación por teclado y ausencia de acciones tapadas por barras fijas.

\clearpage
\section*{Referencias}
\addcontentsline{toc}{section}{Referencias}
\begingroup
\setlength{\parindent}{0pt}
\setlength{\parskip}{0pt}
\everypar{\hangindent=0.5in\hangafter=1}
American Psychological Association. (s. f.). \textit{Paper format}. \url{https://apastyle.apa.org/style-grammar-guidelines/paper-format}\par
xmenil. (2026). \textit{Frontend-Quickly} (commit d0a8c92) [Código fuente]. GitHub. \url{https://github.com/xmenil/Frontend-Quickly/tree/d0a8c92cab013decde2fb45ea01e4daae964389c}\par
\endgroup

\appendix
\clearpage
\section*{Apéndice A}
\addcontentsline{toc}{section}{Apéndice A}
\begin{center}\textbf{Mapa de archivos de las pantallas}\end{center}

\begingroup\singlespacing
\begin{longtable}{L{3.4cm}L{11.9cm}}
\caption{Mapa de archivos}\\
\toprule
Área & Archivos de referencia \\
\midrule
\endhead
Rutas & \path{src/app/router.tsx} \\
Estructura & \path{src/layouts/PublicLayout.tsx}, \path{src/layouts/CustomerLayout.tsx}, \path{src/layouts/MerchantLayout.tsx}, \path{src/layouts/CourierLayout.tsx}, \path{src/layouts/AdminLayout.tsx} \\
Carrito & \path{src/features/cart/CartPage.tsx} \\
Checkout & \path{src/features/checkout/CheckoutPage.tsx}, \path{src/features/checkout/checkoutFlow.ts} \\
Direcciones & \path{src/components/shared/AddressForm.tsx}, \path{src/features/customer/CustomerAddressesPage.tsx} \\
Resumen de precios & \path{src/components/shared/PriceSummary.tsx} \\
Pedidos del cliente & \path{src/features/orders/CustomerOrdersPage.tsx}, \path{src/features/orders/OrderDetailPage.tsx}, \path{src/features/orders/TrackingPage.tsx} \\
Comercio & Pantallas de \path{src/features/merchant/}; preparación en \path{src/features/merchant/MerchantOrdersPage.tsx}. \\
Solicitudes de reparto & \path{src/features/courier/CourierRequestsPage.tsx} \\
Entrega activa & \path{src/features/courier/CourierActiveDeliveryPage.tsx} \\
Historial de reparto & \path{src/features/courier/CourierHistoryPage.tsx} \\
Ganancias & \path{src/features/courier/CourierEarningsPage.tsx} \\
Perfil de reparto & \path{src/features/courier/CourierProfilePage.tsx} \\
Administración & Pantallas de \path{src/features/admin/}; pedidos en \path{src/features/admin/AdminOrdersPage.tsx} y soporte en \path{src/features/admin/AdminTicketsPage.tsx}. \\
Pruebas & \path{src/test/}; incluye validaciones de Checkout, integración de dirección, selección de carrito y operaciones de reparto. \\
\bottomrule
\end{longtable}
\endgroup
\noindent\textit{Nota.} Elaboración propia a partir del código del proyecto (xmenil, 2026).\par\bigskip

\clearpage
\section*{Apéndice B}
\addcontentsline{toc}{section}{Apéndice B}
\begin{center}\textbf{Compilación de este manual}\end{center}

Este archivo utiliza UTF-8 y no depende de imágenes externas. Requiere paquetes habituales de una distribución LaTeX, incluidos \codigo{newtxtext}, \codigo{setspace}, \codigo{fancyhdr}, \codigo{titlesec}, \codigo{caption} y \codigo{babel}, \codigo{hyperref}, \codigo{tikz}, \codigo{longtable} y \codigo{xurl}. El preámbulo selecciona las fuentes según el motor utilizado.

Para generar el PDF y actualizar las referencias de tablas y figuras, desde la carpeta \codigo{docs}:
\begin{verbatim}
pdflatex manual-funcionamiento-quickly.tex
pdflatex manual-funcionamiento-quickly.tex
\end{verbatim}

También puede compilarse dos veces con LuaLaTeX o XeLaTeX. El documento no requiere ejecutar comandos externos desde LaTeX ni habilitar \codigo{shell-escape}.

\end{document}
